\chapter{Verifica e validazione}
\label{cap:verifica_validazione}
L'ultima fase del periodo di tirocinio ha riguardato il processo di verifica e validazione. La prima è un processo che si occupa di fornire evidenza oggettiva, ovvero che i risultati ottenuti come output dello sviluppo del software soddisfino i requisiti, la seconda invece è un processo per confermare in modo definitivo che le caratteristiche del software siano conformi ai bisogni dell'utente e all'uso previsto.
Per eseguire la verifica e la validazione è necessario l'esecuzione di test\footnote{In ambito informatico i test eseguiti nel processo di verifica sono eseguiti in maniera automatica durante lo sviluppo del software} che individuino carenze, correttezza e affidabilità. Nel progetto della Skill Concierge Croccante non è stato realizzato alcun test per il processo di verifica: viste le tempistiche non è stato potuto fare ciò e si è pertanto scelto di eseguire test già esistenti nel servizio Amazon Lambda, sufficienti ad esaminare il codice e rilevare eventuali errori. Per quanto riguarda il processo di validazione\footnote{Attività di supporto la quale accerta che il prodotto dei processi rispetti le specifiche} è stato svolto, alla fine del periodo di tirocinio, un incontro  con il tutor aziendale e CFO Damiano Buscemi come conferma fiale la quale ha accertato le bontà richieste dal prodotto finale. Tale validazione non è stata altro che una dimostrazione funzionante della Skill, eseguita prima da me per spiegare le varie funzionalità, ed infine da Crispy Bacon come conferma.